\documentclass{galahad}

% set the package name

\newcommand{\package}{snls}
\newcommand{\packagename}{SNLS}
\newcommand{\fullpackagename}{\libraryname\_\packagename}
\newcommand{\solver}{{\tt \fullpackagename\_solve}}

\begin{document}

\galheader

%%%%%%%%%%%%%%%%%%%%%% SUMMARY %%%%%%%%%%%%%%%%%%%%%%%%

\galsummary
This package uses a {\bf regularization method to find a (local)
minimizer of a differentiable weighted sum-of-squares objective function}
\eqn{f}{f(\bmx) \eqdef
 \half \sum_{i=1}^{m_r} w_i^{} r_i^2(\bmx) \equiv \half \|\bmr(\bmx)\|_W^2,}
where the variables $\bmx$ are required to lie within the
{\bf regular simplex}
\disp{\bme^T \bmx = 1 \tim{and} \bmx \geq 0,}
or an intersection of {\bf multiple non-overlapping regular simplices}
\eqn{c}{\bme_{\calC_i}^T \bmx_{\calC_i} = 1 \tim{and} \bmx_{\calC_i}
 \geq 0 \tim{for} i = 1,\ldots,m_c,}
of many variables $\bmx$,
the positive weights $w_i$, $i=1,\ldots,m_c$ are given, and
the vector $\bmv_{\calC}$ is made up of those entries of $\bmv$ indexed by
the set $\calC$, 
the index sets of {\bf cohorts} $\calC_i \subseteq \{1,\ldots,n\}$
for which $\calC_i \cap \calC_j = \emptyset$ for $1 \leq i, \neq j \leq m_c$,
The method offers the choice of projected-gradient and interior-point
iteration to solve the key regularization subproblems,
 and is most suitable for large problems.
First derivatives of the {\bf residual function}
$\bmr(\bmx)$ must be available, either directly as the Jacobian matrix 
$\bmJ_r(\bmx) = \nabla \bmr(\bmx)$ or via the action of this matrix
(and its transpose) on specified vectors.
%if second derivatives of the
%$r_i(\bmx)$ can be calculated, they may be exploited---if suitable products
%of the first or second derivatives with a vector may be found but not the
%derivatives themselves, that can also be used to advantage.

%%%%%%%%%%%%%%%%%%%%%% attributes %%%%%%%%%%%%%%%%%%%%%%%%

\galattributes
\galversions{\tt  \fullpackagename\_single, \fullpackagename\_double}.
\galuses
{\tt GALAHAD\_CLOCK},
{\tt GALAHAD\_SY\-M\-BOLS},
{\tt GALAHAD\_NLPT},
{\tt GALAHAD\_USERDATA},
{\tt GALAHAD\_REVERSE},
{\tt GALAHAD\_SPECFILE},
{\tt GALAHAD\_QPT},
{\tt GALAHAD\_SLLS},
{\tt GALAHAD\_SLLSB},
{\tt GALAHAD\_\-SPACE},
{\tt GALAHAD\_MOP},
{\tt GALAHAD\_NORMS}
and
{\tt GALAHAD\_STRING}.
\galdate March 2026.
\galorigin N. I. M. Gould, Rutherford Appleton Laboratory.
\gallanguage Fortran~95 + TR 15581 or Fortran~2003.

%%%%%%%%%%%%%%%%%%%%%% HOW TO USE %%%%%%%%%%%%%%%%%%%%%%%%

\galhowto

\input{versions}

\noindent
If it is required to use more than one of the modules at the same time, 
the derived types
{\tt SMT\_type},
{\tt USERDATA\_type},
{\tt REVERSE\_type},
{\tt \packagename\_time\_\-type},
{\tt \packagename\_control\_type},
{\tt \packagename\_inform\_type},
{\tt \packagename\_data\_type}
and
{\tt NLPT\_problem\_\-type},
(Section~\ref{galtypes})
and the subroutines
{\tt \packagename\_initialize},
{\tt \packagename\_\-solve},
{\tt \packagename\_terminate},
(Section~\ref{galarguments})
and
{\tt \packagename\_\-read\_specfile}
(Section~\ref{galfeatures})
must be renamed on one of the {\tt USE} statements.

%%%%%%%%%%%%%%%%%%%%%% terminology %%%%%%%%%%%%%%%%%%%%%%%%

\galterminology
The algorithm used is iterative. From the current best estimate
of the minimizer $\bmx_k$, a trial improved point $\bmx_k + \bms_k$ is sought.
The correction $\bms_k$ is chosen to improve a model $m_k(\bms)$ of
the objective function $f(\bmx_k+\bms)$, built around $\bmx_k$, within the 
simplex constraint set $\calC(\bmx_k+\bms)$ for which 
\disp{\calC(\bmx) \eqdef \{ \bmx \;|\; \mbox{\req{c} holds}\}.}
The model is the sum of two basic components,
a suitable approximation $t_k(\bms)$ of $f(\bmx_k+\bms)$,
and a regularization term $\sfrac{\sigma_k}{p} \|\bms\|^p$
involving a weight $\sigma_k$ and a power $p$
%and a norm $\|\bms\|_{\bmS_k} \eqdef \sqrt{\bms^T \bmS_k \bms}$ 
%for a given positive definite scaling matrix $\bmS_k$ 
that is included to prevent large corrections.
The weight  $\sigma_k$ is adjusted as the algorithm progresses to
ensure convergence.

The model $t_k(\bms)$ is a truncated Taylor-series approximation, and this
relies on being able to compute or estimate derivatives of $\bmr(\bmx)$.
Various models are provided, and each has different derivative requirements.
We denote the $m_r$ by $n$ {\bf residual/observation Jacobian} 
$\bmJ_r(\bmx)$ as the matrix whose $i,j$-th component
\disp{\bmJ_r(\bmx)_{i,j} \eqdef \partial r_i(\bmx) / \partial x_j \;\;
\mbox{for $i=1,\ldots,o$ and $j=1,\ldots,n$.}}
For a given $m_r$-vector $\bmy$, the
{\bf weighted residual Hessian} is the sum
\disp{\bmH(\bmx,\bmy) \eqdef \sum_{\ell=1}^m y_\ell \bmH_\ell(\bmx),
\;\; \mbox{where}\;\;
\bmH_\ell(\bmx)_{i,j} \eqdef \partial^2 r_\ell(\bmx) / \partial x_i \partial x_j
\;\; \mbox{for $i,j=1,\ldots,n$}}
is the Hessian of $r_\ell(\bmx)$.
The models $t_k(\bms)$ provided are,
\begin{enumerate}
\item the Gauss-Newton approximation
  $\half \| \bmr(\bmx_k) + \bmJ_r(\bmx_k) \bms\|_{\bmW}^2$,
\item the Newton (second-order Taylor) approximation
  $f(\bmx_k) + \bmg(\bmx_k)^T s + \half \bms^T
  [ \bmJ_r^T(\bmx_k) \bmW \bmJ_r(\bmx_k) + \bmH(\bmx_k,\bmW \bmr(\bmx_k))] \bms$,
\end{enumerate}
where $\bmW$ is the diagonal matrix of weights $w_i$, $i = 1, \ldots m$.

\noindent
Access to a particular model requires that the user is either able to
provide the derivatives needed ({\bf ``matrix available''}) or that the products
of these derivatives (and their transposes) with specified vectors are possible
({\bf ``matrix free''}).

The method makes considerable use of efficient available methods for computing 
the projection $P[\bmx]$, that is defined to be the closest (two-norm) point 
of a given $\bmx$ to the simplex constraint set $\calC(\bmx)$.

%%%%%%%%%%%%%%%%%%%%%% matrix storage formats %%%%%%%%%%%%%%%%%%%%%%%%

\galmatrix The matrices $\bmJ_r(\bmx)$ and $\bmH(\bmx,\bmy)$
(as required and when available) may be stored in a variety of input formats.

\subsubsection{Dense storage format}\label{dense}
The matrix $\bmJ_r$ is stored as a compact
dense matrix by rows, that is, the values of the entries of each row in turn are
stored in order within an appropriate real one-dimensional array.
Component $n \ast (i-1) + j$ of the storage array {\tt  Jr\%val} will hold 
the value ${\bmJ_r}_{i,j}$ for $i = 1, \ldots , o$, $j = 1, \ldots , n$. 
Since $\bmH$ is symmetric, only the lower triangular
part (that is the part $\bmH_{ij}$ for $1 \leq j \leq i \leq n$)
should be stored. In this case the lower triangle will be stored by rows,
that is component $i \ast (i-1)/2 + j$ of the storage array {\tt H\%val} will
hold the value $\bmH_{ij}$ (and, by symmetry, $\bmH_{ji}$) for $1 \leq j
\leq i \leq n$.

\subsubsection{Sparse co-ordinate storage format}\label{coordinate}
Only the nonzero entries of the matrices are stored. For the $l$-th
entry of $\bmJ_r$, its row index $i$, column index $j$ and value
${\bmJ_r}_{i,j}$ are stored in the $l$-th components of the integer arrays
{\tt Jr\%row}, {\tt Jr\%col} and real array {\tt Jr\%val}.
The order is unimportant, but the total number of entries {\tt Jr\%ne} is
required.  The same scheme is applicable to $\bmH$
(thus requiring integer arrays {\tt H\%row}, {\tt H\%col}, a real
array {\tt H\%val}, and an integer value {\tt H\%ne}), except that only
the entries in the lower triangle should be stored.

%We recommend this storage format since it is the most efficient
%storage scheme for the underlying linear system solvers.  Any other
%storage formats are converted internally to sparse co-ordinate
%storage and then handled accordingly.

\subsubsection{Sparse row-wise storage format}\label{rowwise}
Again only the nonzero entries are stored, but this time they are
ordered so that those in row $i$ appear directly before those in row
$i+1$. For the $i$-th row of $\bmJ_r$, the $i$-th component of the integer
array {\tt Jr\%ptr} holds the position of the first entry in this row,
while {\tt Jr\%ptr} $(m_r+1)$ holds the total number of entries plus one.
The column indices $j$ and values ${\bmJ_r}_{i,j}$ of the entries in the
$i$-th row are stored in components $l =$ {\tt Jr\%ptr}$(i)$, \ldots
,{\tt Jr\%ptr} $(i+1)-1$ of the integer array {\tt Jr\%col}, and real
array {\tt Jr\%val}, respectively.
The same scheme is applicable to
$\bmH$ (thus requiring integer arrays {\tt H\%ptr}, {\tt H\%col}, and
a real array {\tt H\%val}),
except that only the entries in the lower triangle should be stored.

For sparse matrices, this scheme almost always requires less storage than
its predecessor.

\subsubsection{Sparse column-wise storage format}\label{columnwise}
For the matrix $\bmJ_r$, once again only the nonzero entries are
stored, but this time they are ordered so that those in column $j$ appear
directly before those in column $j+1$. For the $j$-th column of $\bmJ_r$, the
$j$-th component of the integer array {\tt Jr\%ptr} holds the position of the
first entry in this column, while {\tt Jr\%ptr} $(n+1)$ holds the total number
of entries plus one.  The row indices $i$ and values ${\bmJ_r}_{ij}$ of the
entries in the $j$-th column are stored in components $l =$ 
{\tt Jr\%ptr}$(j)$, \ldots ,{\tt Jr\%ptr} $(j+1)-1$ of the integer array
{\tt Jr\%row}, and real array {\tt Jr\%val}, respectively.

\subsubsection{Diagonal storage format}\label{diagonal}
If $\bmH$ is diagonal (i.e., $\bmH_{ij} = 0$ for all $1 \leq i \neq j \leq n$)
only the diagonal entries $\bmH_{ii}$ for $1 \leq i \leq n$ need be stored,
and the first $n$ components of the array {\tt H\%val} may be used for
the purpose. There is no sensible equivalent for the non-square
matrix $\bmJ_r$.

%%%%%%%%%%%%%%%%%%%%%%%%%%% kinds %%%%%%%%%%%%%%%%%%%%%%%%%%%

\input{kinds}

%%%%%%%%%%%%%%%%%%%%%% derived types %%%%%%%%%%%%%%%%%%%%%%%%

\galtypes
Eight derived data types are accessible from the package.

%%%%%%%%%%% matrix data type %%%%%%%%%%%

\subsubsection{The derived data type for holding matrices}\label{typesmt}
The derived data type {\tt SMT\_TYPE} is used to hold the Hessian matrix $\bmH$
if this is available. The components of {\tt SMT\_TYPE} used here are:

\begin{description}

\ittf{m} is a scalar component of type \integer,
that holds the row dimension of the matrix.

\ittf{n} is a scalar component of type \integer,
that holds the column dimension of the matrix.

\ittf{ne} is a scalar variable of type \integer, that
holds the number of matrix entries.

\ittf{type} is a rank-one allocatable array of type default \character, that
is used to indicate the matrix storage scheme used. Its precise length and
content depends on the type of matrix to be stored (see \S\ref{typeprob}).

\ittf{val} is a rank-one allocatable array of type \realdp\,
and dimension at least {\tt ne}, that holds the values of the entries.
Each pair of off-diagonal entries $h_{ij} = h_{ji}$ of the {\em symmetric}
matrix $\bmH$ is represented as a single entry
(see \S\ref{dense}--\ref{diagonal}).
Any duplicated entries that appear in the sparse
co-ordinate or row-wise schemes will be summed.

\ittf{row} is a rank-one allocatable array of type \integer,
and dimension at least {\tt ne}, that may hold the row indices of the entries.
(see \S\ref{coordinate} and \S\ref{columnwise}).

\ittf{col} is a rank-one allocatable array of type \integer,
and dimension at least {\tt ne}, that may hold the column indices of the entries
(see \S\ref{coordinate}--\ref{rowwise}).

\ittf{ptr} is a rank-one allocatable array of type \integer,
and dimension at least {\tt m + 1}, that may hold the pointers to
the first entry in each row (see \S\ref{rowwise}--\ref{columnwise}).

\end{description}

%%%%%%%%%%% problem type %%%%%%%%%%%

\subsubsection{The derived data type for holding the problem}\label{typeprob}
The derived data type {\tt NLPT\_problem\_type} is used to hold
the problem. The relevant components of
{\tt NLPT\_problem\_type}
are:

\begin{description}

\ittf{n} is a scalar variable of type \integer,
 that holds the number of optimization variables, $n$.

\ittf{m\_r} is a scalar variable of type \integer,
 that holds the number of residual functions, $m_r$.

\ittf{m\_c} is a scalar variable of type \integer,
 that holds the number of cohorts, $m_c$.

%\itt{stabilisation\_weight}  is a scalar variable of type \realdp,
%that holds the value of the stabilisation weight, $\rho \geq 0$.

%\itt{stabilisation\_power}  is a scalar variable of type \realdp,
%that holds the value of the stabilisation power, $r \geq 2$.

\ittf{Jr} is scalar variable of type {\tt SMT\_TYPE}
that holds the Jacobian matrix $\bmJ_r(\bmx)$ (if it is available).
The following components are used here:

\begin{description}

\itt{Jr\%type} is an allocatable array of rank one and type default
\character, that
is used to indicate the storage scheme used. If the dense storage scheme
(see Section~\ref{dense}) is used,
the first five components of {\tt Jr\%type} must contain the
string {\tt DENSE}.
For the sparse co-ordinate scheme (see Section~\ref{coordinate}),
the first ten components of {\tt Jr\%type} must contain the
string {\tt COORDINATE},
and for the sparse row-wise storage scheme (see Section~\ref{rowwise}),
the first fourteen components of {\tt Jr\%type} must contain the
string {\tt SPARSE\_BY\_ROWS}.

For convenience, the procedure {\tt SMT\_put}
may be used to allocate sufficient space and insert the required keyword
into {\tt Jr\%type}.
For example, if {\tt nlp} is of derived type {\tt \packagename\_problem\_type}
and involves a Jacobian we wish to store using the co-ordinate scheme,
we may simply
%\vspace*{-2mm}
{\tt
\begin{verbatim}
        CALL SMT_put( nlp%J%type, 'COORDINATE' )
\end{verbatim}
}
%\vspace*{-4mm}
\noindent
See the documentation for the \galahad\ package {\tt SMT}
for further details on the use of {\tt SMT\_put}.

\itt{Jr\%ne} is a scalar variable of type \integer, that
holds the number of entries in $\bmJ_r$
in the sparse co-ordinate storage scheme (see Section~\ref{coordinate}).
It need not be set for any of the other two schemes.

\itt{Jr\%val} is a rank-one allocatable array of type \realdp, that holds
the values of the entries of the Jacobian matrix $\bmJ_r$ in any of the
storage schemes discussed in Section~\ref{galmatrix}.

\itt{Jr\%row} is a rank-one allocatable array of type \integer,
that holds the row indices of $\bmJ_r$ in the sparse co-ordinate storage
scheme (see Section~\ref{coordinate}).
It need not be allocated for any of the other two schemes.

\itt{Jr\%col} is a rank-one allocatable array variable of type \integer,
that holds the column indices of $\bmJ_r$ in either the sparse co-ordinate
(see Section~\ref{coordinate}), or the sparse row-wise
(see Section~\ref{rowwise}) storage scheme.
It need not be allocated when the dense storage scheme is used.

\itt{Jr\%ptr} is a rank-one allocatable array of dimension {\tt m+1} and type
\integer, that holds the starting position of each row of $\bmJ_r$, as well
as the total number of entries plus one, in the sparse row-wise storage
scheme (see Section~\ref{rowwise}). It need not be allocated when the
other schemes are used.

\end{description}

\ittf{H} is scalar variable of type {\tt SMT\_TYPE} that may hold the
weighted Hessian matrix $\bmH(\bmx,\bmy)$ (if it is available).
The following components are used here:

\begin{description}

\itt{H\%type} is an allocatable array of rank one and type default
\character, that
is used to indicate the storage scheme used. If the dense storage scheme
(see Section~\ref{dense}) is used,
the first five components of {\tt H\%type} must contain the
string {\tt DENSE}.
For the sparse co-ordinate scheme (see Section~\ref{coordinate}),
the first ten components of {\tt H\%type} must contain the
string {\tt COORDINATE},
for the sparse row-wise storage scheme (see Section~\ref{rowwise}),
the first fourteen components of {\tt H\%type} must contain the
string {\tt SPARSE\_BY\_ROWS},
and for the diagonal storage scheme (see Section~\ref{diagonal}),
the first eight components of {\tt H\%type} must contain the
string {\tt DIAGONAL}.

For convenience, the procedure {\tt SMT\_put}
may be used to allocate sufficient space and insert the required keyword
into {\tt H\%type}.
For example, if {\tt nlp} is of derived type {\tt \packagename\_problem\_type}
and involves a Hessian we wish to store using the co-ordinate scheme,
we may simply
%\vspace*{-2mm}
{\tt
\begin{verbatim}
        CALL SMT_put( nlp%H%type, 'COORDINATE' )
\end{verbatim}
}
%\vspace*{-4mm}
\noindent
See the documentation for the \galahad\ package {\tt SMT}
for further details on the use of {\tt SMT\_put}.

\itt{H\%ne} is a scalar variable of type \integer, that
holds the number of entries in the {\bf lower triangular} part of $\bmH$
in the sparse co-ordinate storage scheme (see Section~\ref{coordinate}).
It need not be set for any of the other three schemes.

\itt{H\%val} is a rank-one allocatable array of type \realdp, that holds
the values of the entries of the {\bf lower triangular} part
of the Hessian matrix $\bmH$ in any of the
storage schemes discussed in Section~\ref{galmatrix}.

\itt{H\%row} is a rank-one allocatable array of type \integer,
that holds the row indices of the {\bf lower triangular} part of $\bmH$
in the sparse co-ordinate storage
scheme (see Section~\ref{coordinate}).
It need not be allocated for any of the other three schemes.

\itt{H\%col} is a rank-one allocatable array variable of type \integer,
that holds the column indices of the {\bf lower triangular} part of
$\bmH$ in either the sparse co-ordinate
(see Section~\ref{coordinate}), or the sparse row-wise
(see Section~\ref{rowwise}) storage scheme.
It need not be allocated when the dense or diagonal storage schemes are used.

\itt{H\%ptr} is a rank-one allocatable array of dimension {\tt n+1} and type
\integer, that holds the starting position of
each row of the {\bf lower triangular} part of $\bmH$, as well
as the total number of entries plus one, in the sparse row-wise storage
scheme (see Section~\ref{rowwise}). It need not be allocated when the
other schemes are used.

\end{description}

\itt{COHORT} is a rank-one allocatable array of dimension {\tt n} and type
\integer, that specifies with which cohort variable $x_j$, $j = 1,\ldots,n$, 
is associated. If variable $x_j$ is associated with cohort $\calC_i$, 
$1 \leq i \leq m$, {\tt COHORT(j)} should be set to {\tt i}, while if $x_j$ 
is unconstrained {\tt COHORT(j) = 0} should be assigned. At least one value 
{\tt COHORT(j)} for $j = 1,\ldots\,n$ is expected to take the value $i$ 
for every $1 \leq i \leq m$, that is no empty cohorts are allowed.
If {\tt COHORT} is not allocated, a single cohort containing all of
the variables will be used instead, i.e., a single unit simplex is employed.

\ittf{X} is a rank-one allocatable array of dimension {\tt n} and type
\realdp, that holds the values $\bmx$ of the optimization variables.
The $j$-th component of {\tt X}, $j = 1, \ldots, n$, contains $x_{j}$.

\ittf{Y} is a rank-one allocatable array of dimension {\tt m\_c} (or {\tt 1} if
a single unit simplex is implied) and type default \realdp, that holds
the values $\bmy$ of estimates of the Lagrange multipliers
corresponding to the equality constraint for each cohort 
(see Section~\ref{galmethod}).
The $i$-th component of {\tt Y}, $i = 1, \ldots, m_c$, contains $y_{i}$.

\ittf{Z} is a rank-one allocatable array of dimension {\tt n} and type default
\realdp, that holds the values $\bmz$ of estimates  of the dual variables
corresponding to the non-negativity constraints (see Section~\ref{galmethod}).
The $j$-th component of {\tt Z}, $j = 1, \ldots, n$, contains $z_{j}$.

\ittf{R} is a rank-one allocatable array of dimension {\tt o} and type
\realdp, that holds the residual values $\bmr(\bmx)$
at the point $\bmx$. The $i$-th component of
{\tt R}, $i = 1, \ldots, o$, contains $\bmr_{i}(\bmx)$.

\ittf{G} is a rank-one allocatable array of dimension {\tt m} and type
\realdp, that holds the gradient $g(\bmx) = \nabla_x f(\bmx)$ of the
objective function at the point $\bmx$. The $j$-th component of
{\tt G}, $i = 1, \ldots, n$, contains $\partial f(\bmx)/\partial x_j$.

\itt{X\_status} is a rank-one allocatable array argument of dimension {\tt n}
and type \integer, that indicates which of the non-negativity constraints
are in the current active set. Possible values for
{\tt X\_status(j)}, {\tt j}$=1, \ldots ,$ {\tt n}, and their meanings are
\begin{description}
\itt{<0} the $j$-th non-negativity constraint
is in the active set, at the value zero, and
\itt{0}  the $j$-th non-negativity constraint is not in the active set,
or the associated variable is unconstrained.
\end{description}

\end{description}

%%%%%%%%%%% control type %%%%%%%%%%%

\subsubsection{The derived data type for holding control
 parameters}\label{typecontrol}
The derived data type
{\tt \packagename\_control\_type}
is used to hold controlling data. Default values may be obtained by calling
{\tt \packagename\_initialize}
(see Section~\ref{subinit}),
while components may also be changed by calling
{\tt \fullpackagename\_read\-\_spec}
(see Section~\ref{readspec}).
The derived type {\tt \packagename\_subproblem\_control\_type}
comprises all of the components of {\tt \packagename\_control\_type}
except for the last (i.e., {\tt subproblem\_control}).
The components of
{\tt \packagename\_control\_type}
are:

\begin{description}

\itt{error} is a scalar variable of type \integer, that holds the
stream number for error messages. Printing of error messages in
{\tt \packagename\_solve} and {\tt \packagename\_terminate}
is suppressed if {\tt error} $\leq 0$.
The default is {\tt error = 6}.

\ittf{out} is a scalar variable of type \integer, that holds the
stream number for informational messages. Printing of informational messages in
{\tt \packagename\_solve} is suppressed if {\tt out} $< 0$.
The default is {\tt out = 6}.

\itt{print\_level} is a scalar variable of type \integer, that is used
to control the amount of informational output which is required. No
informational output will occur if {\tt print\_level} $\leq 0$. If
{\tt print\_level} $= 1$, a single line of output will be produced for each
iteration of the process. If {\tt print\_level} $\geq 2$, this output will be
increased to provide significant detail of each iteration.
The default is {\tt print\_level = 0}.

\itt{start\_print} is a scalar variable of type \integer, that specifies
the first iteration for which printing will occur in {\tt \packagename\_solve}.
If {\tt start\_print} is negative, printing will occur from the outset.
The default is {\tt start\_print = -1}.

\itt{stop\_print} is a scalar variable of type \integer, that specifies
the last iteration for which printing will occur in  {\tt \packagename\_solve}.
If {\tt stop\_print} is negative, printing will occur once it has been
started by {\tt start\_print}.
The default is {\tt stop\_print = -1}.

\itt{print\_gap} is a scalar variable of type \integer.
Once printing has been started, output will occur once every
{\tt print\_gap} iterations. If {\tt print\_gap} is no larger than 1,
printing will be permitted on every iteration.
The default is {\tt print\_gap = 1}.

\itt{maxit} is a scalar variable of type \integer, that holds the
maximum number of iterations which will be allowed in {\tt \packagename\_solve}.
The default is {\tt maxit = 1000}.

\itt{alive\_unit} is a scalar variable of type \integer.
If {\tt alive\_unit} $>$ 0, a temporary file named {\tt alive\_file} (see below)
will be created on stream number {\tt alive\_unit} on initial entry to
\solver, and execution of \solver\ will continue so
long as this file continues to exist. Thus, a user may terminate execution
simply by removing the temporary file from this unit.
If {\tt alive\_unit} $\leq$ 0, no temporary file will be created, and
execution cannot be terminated in this way.
The default is {\tt alive\_unit} $=$ 40.

\itt{jacobian\_available} is a scalar variable of type \integer,
that specifies the availability of the residual Jacobian.
Possible values are:

\begin{description}
\itt{$\geq$ 2} The Jacobian is available.
\itt{1} The Jacobian is not available explicitly but its effect
may be accessed by matrix-vector products, i.e., it is ``matrix-free''.
\itt{$\leq$ 0} The Jacobian is not available either explicitly or via
matrix-vector products.
\end{description}
The default is {\tt jacobian\_available = 1}.

\itt{subproblem\_solver} is a scalar variable of type \integer, that
specifies the method used to solve the crucial step-determination
subproblem. Possible values are:
\begin{description}
\itt{$\leq$ 1} A projected-gradient method using {\tt GALAHAD\_SLLS} 
will be used.
\itt{2} An interior-point method using {\tt GALAHAD\_SLLSB} will be used.
\itt{$\geq$ 3} An interior-point method will initially be used,  but a
switch to a projected-gradient method subsequently occur 
(see {\tt \%stop\_pg\_switch}.
\end{description}
The default is {\tt subproblem\_solver = 1}.

\itt{non\_monotone} is a scalar variable of type \integer, that
specifies the history-length for non-monotone descent strategy.
Any non-positive value results in standard monotone descent, for which
merit function improvement occurs at each iteration. There are often
definite advantages in using a non-monotone strategy with a modest history,
since the occasional local increase in the merit function may
enable the algorithm to move across (gentle) ``ripples'' in
the merit function surface.
However, we do not usually recommend large values of {\tt non\_monotone}.
The default is {\tt non\_monotone = 1}.

\itt{weight\_update\_strategy} is a scalar variable of type \integer,
that specifies the way in which the regularization weight will be adjusted
at the end of each iteration. Possible values are:
\begin{description}
\itt{1} the traditional acceptance and rejection strategy as described below.
\itt{2} the traditional strategy, except that a zero weight will be tried
first after a very successful step.
\itt{3} a more sophisticated strategy that mimics that proposed
for trust-region methods by Gould, Porcelli and Toint.
\end{description}
Any other value will be considered as if {\tt weight\_update\_strategy = 1},
and this is the default.

\itt{stop\_r\_absolute} is a scalar variable of type \realdp,
that is used to specify the maximum permitted norm of the
residual, $\|\bmr(\bmx)\|_{\bmW}$,
(see Section~\ref{galmethod}) at the estimate of the solution sought.
Any negative value will be recorded as zero.
The default is {\tt stop\_r\_absolute =} $10^{-6}$.

\itt{stop\_r\_relative} is a scalar variable of type \realdp,
that is used to specify the largest relative reduction in the two-norm of the
residual, $\|\bmr(\bmx)\|_{\bmW}$, that will be permitted
(see Section~\ref{galmethod}) at the estimate of the solution sought
compared to that at the initial point.
Any negative value will be recorded as zero.
The default is {\tt stop\_r\_relative =} $0.0$.

\itt{stop\_pg\_absolute} is a scalar variable of type \realdp, that is used 
to specify the maximum permitted norm of the projected gradient,
$\|P[\bmx - \bmJ_r^T(\bmx) \bmW \bmr(\bmx)] - \bmx\|_2$,
of the residual  $\|\bmr(\bmx)\|_{\bmW}$,
(see Section~\ref{galmethod}) at the estimate of the solution sought.
Any negative value will be recorded as zero.
The default is {\tt stop\_pg\_absolute =} $10^{-6}$.

\itt{stop\_pg\_relative} is a scalar variable of type \realdp,
that is used to specify the largest relative reduction in the norm of the
projected gradient of the residual that will be permitted
(see Section~\ref{galmethod}) at the estimate of the solution sought
compared to that at the initial point.
Any negative value will be recorded as zero.
The default is {\tt stop\_pg\_relative =} $0.0$.
%The default is {\tt stop\_pg\_relative =} $10^{-6}$.

\itt{stop\_s} is a scalar variable of type \realdp,
that is used to specify the minimum acceptable correction step $\bms$
relative to the current estimate of the solution $\bmx$
The algorithm will be deemed to have converged if $|s_i| \leq$
{\tt stop\_s} $\ast \max( 1, |x_i|)$ for all $i = 1, \ldots, n$.
The default is {\tt stop\_s =} $u$,
where $u$ is {\tt EPSILON(1.0)} ({\tt EPSILON(1.0D0)} in
{\tt \fullpackagename\_double}).

\itt{stop\_pg\_switch} is a scalar variable of type \realdp, that is used 
to specify the value of the norm of the projected gradient below which
the subproblem solver will change from an interior-point to a projection 
method when {\tt \%subproblem\_solver = 3}. 
The default is {\tt stop\_pg\_switch = 0.001}.

\itt{initial\_weight} is a scalar variable of type \realdp, that holds
the required initial value of the regularization weight. If
{\tt initial\_weight} $\leq 0$, the weight will be chosen automatically
by \solver.
%The default is {\tt initial\_weight = - 1.0}.
The default is {\tt initial\_weight = 100.0}.

\itt{minimum\_weight} is a scalar variable of type \realdp, that holds
the largest permitted value of the regularization weight as the algorithm
proceeds.
%If  {\tt minimum\_weight} $\leq 0$, the weight will be chosen automatically
%by \solver\
The default is {\tt minimum\_weight =} $10^{-8}$.

\ittt{eta\_successful}, {\tt eta\_very\_successful}
and {\tt eta\_too\_successful}
are scalar variables of type default
\realdp, that control the acceptance and rejection of the trial step
and the updates to the regularization weight.
At every iteration, the ratio of the actual reduction in the merit function
following the trial step to that predicted by the model is computed.
The step is accepted whenever this ratio exceeds {\tt eta\_successful};
otherwise the regularization weight will be reduced.
If, in addition, the ratio exceeds {\tt eta\_very\_successful} but not
{\tt eta\_too\_successful}, the regularization weight may be increased.
The defaults are
{\tt eta\_successful =} $10^{-8}$,
{\tt eta\_very\_successful = 0.9} and
{\tt eta\_too\_successful = 2.0}.

\ittt{weight\_increase},
{\tt weight\_decrease},
{\tt weight\_increase\_max}
and
{\tt weight\_decrease\_min}
are scalar variables of type \realdp, that
control the maximum amounts by which the regularization weight can
contract or expand during an iteration. The weight will be decreased by
powers of {\tt weight\_decrease}, but not in total
more than  {\tt weight\_decrease\_min}, until it is smaller than the norm of
the current step.
It can be increased by at most a factor {\tt weight\_increase}, but not
in total less than {\tt weight\_increase\_max}.
The defaults are
{\tt weight\_increase = 10.0},
{\tt weight\_decrease = 0.1},
{\tt weight\_incr\-ease\_max = 100.0} and
{\tt weight\_decrease\_min = 0.1}.

\itt{switch\_to\_newton} is a scalar variable of type \realdp,
that is used to specify the value of the two-norm of the projected gradient,
required before a switch is made from the Gauss-Newton model to the Newton 
one when {\tt \%newton\_acceleration} is \true.
The default is %{\tt switch\_to\_newton = 0.1}. 
{\bf Not yet implemented.}

\itt{cpu\_time\_limit} is a scalar variable of type \realdp,
that is used to specify the maximum permitted CPU time. Any negative
value indicates no limit will be imposed. The default is
{\tt cpu\_time\_limit = - 1.0}.

\itt{clock\_time\_limit} is a scalar variable of type \realdp,
that is used to specify the maximum permitted elapsed system clock time.
Any negative value indicates no limit will be imposed. The default is
{\tt clock\_time\_limit = - 1.0}.

\itt{newton\_acceleration}
is a scalar variable of type default \logical,
that should be set \true\ if available  second derivatives should 
be used to accelerate the convergence of the algorithm.
The default is {\tt newton\_acceleration = .FALSE.}.
{\bf Not yet implemented.}

\itt{magic\_step}
is a scalar variable of type default \logical,
that should be set \true\ if additional ``magic'' steps are to be
used in order to improve the objective as the iteration proceeds,
and \false\ if no ``magic'' steps are used.
The default is {\tt magic\_step = .FALSE.}.

\itt{print\_obj}
is a scalar variable of type default \logical,
that should be set \true\ if output values relate to $f(\bmx)$,
and \false\ if they relate to $\|\bmr(\bmx)\|_{\bmW}$.
The default is {\tt print\_obj = .FALSE.}.

\itt{space\_critical} is a scalar variable of type default \logical,
that must be set \true\ if space is critical when allocating arrays
and  \false\ otherwise. The package may run faster if
{\tt space\_critical} is \false\ but at the possible expense of a larger
storage requirement. The default is {\tt space\_critical = .FALSE.}.

\itt{deallocate\_error\_fatal} is a scalar variable of type default \logical,
that must be set \true\ if the user wishes to terminate execution if
a deallocation  fails, and \false\ if an attempt to continue
will be made. The default is {\tt deallocate\_error\_fatal = .FALSE.}.

\itt{alive\_file} is a scalar variable of type default \character\ and length
30, that gives the name of the temporary file whose removal from stream number
{\tt alive\_unit} terminates execution of \solver.
The default is {\tt alive\_unit} $=$ {\tt ALIVE.d}.

\itt{prefix} is a scalar variable of type default \character\
and length 30, that may be used to provide a user-selected
character string to preface every line of printed output.
Specifically, each line of output will be prefaced by the string
{\tt prefix(2:LEN(TRIM( prefix ))-1)},
thus ignoreing the first and last non-null components of the
supplied string. If the user does not want to preface lines by such
a string, they may use the default {\tt prefix = ""}.

\itt{SLLS\_control} is a scalar variable of type
{\tt SLLS\_control\_type}
whose components are used to control
the projected-gradient step calculation using
{\tt \libraryname\_SLLS} has been selected (see {\tt subproblem\_solver}.
See the specification sheet for the package
{\tt \libraryname\_SLLS}
for details, and appropriate default values.

\itt{SLLSB\_control} is a scalar variable of type
{\tt SLLSB\_control\_type}
whose components are used to control
the interior-point step calculation using
{\tt \libraryname\_SLLB} has been selected (see {\tt \%subproblem\_solver}.
See the specification sheet for the package
{\tt \libraryname\_SLLB}
for details, and appropriate default values.

\end{description}

%%%%%%%%%%% time type %%%%%%%%%%%

\subsubsection{The derived data type for holding timing
 information}\label{typetime}
The derived data type
{\tt \packagename\_time\_type}
is used to hold elapsed CPU and system clock times for the various parts
of the calculation. The components of
{\tt \packagename\_time\_type}
are:
\begin{description}
\itt{total} is a scalar variable of type default \realdp, that gives
 the CPU total time spent in the package.

\itt{slls} is a scalar variable of type \realdp, that gives
the total CPU time spent solving the step-determination subproblem
using the projected-gradient method from the {\tt GALAHAD\_SLLS} package.

\itt{sllsb} is a scalar variable of type \realdp, that gives
the total CPU time spent solving the step-determination subproblem
using the interior-point method from the {\tt GALAHAD\_SLLSB} package.

\itt{clock\_total} is a scalar variable of type default \realdp, that gives
 the total elapsed system clock time spent in the package.

\itt{clock\_slls} is a scalar variable of type \realdp, that gives the total 
elapsed system clock time spent solving the step-determination subproblem
using the projected-gradient method from the {\tt GALAHAD\_SLLS} package.

\itt{clock\_sllsb} is a scalar variable of type \realdp, that gives the total 
elapsed system clock time spent solving the step-determination subproblem
using the interior-point method from the {\tt GALAHAD\_SLLSB} package.

\end{description}

%%%%%%%%%%% inform type %%%%%%%%%%%

\subsubsection{The derived data type for holding informational
 parameters}\label{typeinform}
The derived data type
{\tt \packagename\_inform\_type}
is used to hold parameters that give information about the progress and needs
of the algorithm.
The derived type {\tt \packagename\_subproblem\_inform\_type}
comprises all of the components of {\tt \packagename\_inform\_type}
except for the last (i.e., {\tt subproblem\_inform}).
The components of
{\tt \packagename\_inform\_type}
are:

\begin{description}
\itt{status} is a scalar variable of type \integer, that gives the
exit status of the algorithm.
%See Sections~\ref{galerrors} and \ref{galinfo}
See Sections~\ref{reverse} and \ref{galerrors}
for details.

\itt{alloc\_status} is a scalar variable of type \integer, that gives
the status of the last attempted array allocation or deallocation.
This will be 0 if {\tt status = 0}.

\itt{bad\_alloc} is a scalar variable of type default \character\
and length 80, that  gives the name of the last internal array
for which there were allocation or deallocation errors.
This will be the null string if {\tt status = 0}.

\ittf{iter} is a scalar variable of type \integer, that holds the
number of iterations performed.

\itt{inner\_iter} is a scalar variable of type \integer, that gives the
total number of inner (projected gradient and/or interior-point) iterations 
required.

\itt{c\_eval} is a scalar variable of type \integer, that gives the
total number of residual function evaluations performed.

\itt{Jr\_eval} is a scalar variable of type \integer, that gives the
total number of residual Jacobian evaluations performed.

\ittf{obj} is a scalar variable of type \realdp, that holds the
value of the objective function $f(\bmx)$  at the best estimate
of the solution found.

\itt{norm\_r} is a scalar variable of type \realdp, that holds the
value of the two-norm of the residual function, $\|\bmr(\bmx)\|_{\bmW}$.
at the best estimate of the solution found.

\itt{norm\_g} is a scalar variable of type \realdp, that holds the
value of the norm of the gradient of the two-norm of the residual function,
$\|\bmJ_r^T(\bmx) \bmW \bmr(\bmx)\|_2/\|\bmr(\bmx)\|_{\bmW}$, at the best estimate
of the solution found.

\itt{norm\_pg} is a scalar variable of type \realdp, that holds the
value of the norm of the projected gradient of the residual function,
$\|P[\bmx - \bmJ_r^T(\bmx) \bmW \bmr(\bmx)] - \bmx\|_2$, at the best 
estimate of the solution found.

\itt{weight} is a scalar variable of type \realdp, that holds the
final value of the regularization weight used.

\ittf{time} is a scalar variable of type {\tt \packagename\_time\_type}
whose components are used to hold elapsed elapsed CPU and system clock
times for the various parts of the calculation (see Section~\ref{typetime}).

\itt{SLLS\_inform} is a scalar variable of type
{\tt SLLS\_inform\_type}
whose components give information about the progress and needs
of the projected-gradient subproblem package
{\tt \libraryname\_SLLS}.
See the specification sheet for the package
{\tt \libraryname\_SLLS}
for details.

\itt{SLLSB\_inform} is a scalar variable of type
{\tt SLLSB\_inform\_type}
whose components give information about the progress and needs
of the integer-point subproblem package
{\tt \libraryname\_SLLSB}.
See the specification sheet for the package
{\tt \libraryname\_SLLSB}
for details.

\end{description}

%%%%%%%%%%% data type %%%%%%%%%%%

\subsubsection{The derived data type for holding problem data}\label{typedata}
The derived data type
{\tt \packagename\_data\_type}
is used to hold all the data for a particular problem,
or sequences of problems with the same structure, between calls of
{\tt \packagename} procedures.
This data should be preserved, untouched (except as directed on
return from \solver\ with positive values of {\tt inform\%status}, see
Section~\ref{reverse}),
from the initial call to
{\tt \packagename\_initialize}
to the final call to
{\tt \packagename\_terminate}.

%%%%%%%%%%% userdata type %%%%%%%%%%%

\input{userdata_type}

%%%%%%%%%%% reverse type %%%%%%%%%%%

\subsubsection{The derived data type for holding reverse-communication data}\label{typereverse}
The derived data type
{\tt REVERSE\_type}
is used to hold data needed for reverse communication when this is
required.
The components of
{\tt REVERSE\_type}
are:

\begin{description}

\itt{lvl} is a scalar variable of type \integer, that
may be used to hold the starting position in {\tt V} (see below)
of the list of indices of nonzero components of $\bmv$.

\itt{lvu} is a scalar variable of type \integer, that
may be used to hold the finishing position in {\tt V} (see below)
of the list of indices of nonzero components of $\bmv$.

\itt{IV} is a rank-one allocatable array of dimension $n$
and type \integer, that may be used to hold the indices of the
nonzero components of $\bmv$. If used, components
{\tt IV(lvl:lvu)} of {\tt V} (see below) will be nonzero.

\itt{V} is a rank-one allocatable array of dimension $n$
and type \realdp, that is used to hold the components of an
input vector $\bmv$.

\itt{index} is a scalar variable of type \integer, that
may be used to hold a column index of $\bmJ_r$ if required.

\itt{lp} is a scalar variable of type \integer, that
is used to record the finishing position in 
{\tt P} and {\tt IP} (see below)
of the list of indices of nonzero components a requested operation on $\bmv$.

\itt{P} is a rank-one allocatable array of dimension $m_r$
and type \realdp, that is used to record the components of the
vector that results from a requested operation on $\bmv$.

\itt{IP} is a rank-one allocatable array of dimension $n$
and type \integer, that is used to record the list of
indices of nonzero components of a requested operation on $\bmv$.

\end{description}



%%%%%%%%%%%%%%%%%%%%%% argument lists %%%%%%%%%%%%%%%%%%%%%%%%

\galarguments
There are three procedures for user calls
(see Section~\ref{galfeatures} for further features):

\begin{enumerate}
\item The subroutine
      {\tt \packagename\_initialize}
      is used to set default values, and initialize private data,
      before solving one or more problems with the
      same sparsity and bound structure.
\item The subroutine
      {\tt \packagename\_solve}
      is called to solve the problem.
\item The subroutine
      {\tt \packagename\_terminate}
      is provided to allow the user to automatically deallocate array
       components of the private data, allocated by
       {\tt \packagename\_solve},
       at the end of the solution process.
       It is important to do this if the data object is re-used for another
       problem {\bf with a different structure}
       since {\tt \packagename\_initialize} cannot test for this situation,
       and any existing associated targets will subsequently become unreachable.
\end{enumerate}
We use square brackets {\tt [ ]} to indicate \optional\ arguments.

%%%%%% initialization subroutine %%%%%%

\subsubsection{The initialization subroutine}\label{subinit}
 Default values are provided as follows:
\vspace*{1mm}

\hspace{8mm}
{\tt CALL \packagename\_initialize( data, control, inform )}

\vspace*{-3mm}
\begin{description}

\itt{data} is a scalar \intentinout\ argument of type
{\tt \packagename\_data\_type}
(see Section~\ref{typedata}). It is used to hold data about the problem being
solved.

\itt{control} is a scalar \intentout\ argument of type
{\tt \packagename\_control\_type}
(see Section~\ref{typecontrol}).
On exit, {\tt control} contains default values for the components as
described in Section~\ref{typecontrol}, while
{\tt control\%subproblem\_control} contains default values that are required for the
model minimization if {\tt control\%model} $\geq$ {\tt 6}).
These values should only be changed after calling
{\tt \packagename\_initialize}.

\itt{inform} is a scalar \intentout\ argument of type
{\tt \packagename\_inform\_type}
(see Section~\ref{typeinform}). A successful call to
{\tt \packagename\_initialize}
is indicated when the  component {\tt inform\%status} has the value 0.
For other return values of {\tt inform\%status}, see
Section~\ref{galerrors}.
The components of {\tt inform} correspond to the main algorithm, while
those in {\tt inform\%subproblem\_inform} refer to the model minimization
if {\tt control\%model} $\geq$ {\tt 6}).

\end{description}

%%%%%%%%% main solution subroutine %%%%%%

\subsubsection{The minimization subroutine}
The minimization algorithm is called as follows:
\vspace*{1mm}

\hspace{8mm}
{\tt CALL \packagename\_solve( nlp, control, inform, data, userdata[, 
   reverse, eval\_R, eval\_Jr, \hspace{3mm}                  \&}
\vspace*{-5mm}

\hspace{37mm}
{\tt eval\_Jr\_prod, eval\_Jr\_scol,  eval\_Jr\_sprod ] )}

\vspace*{-2mm}
\begin{description}
\ittf{nlp} is a scalar \intentinout\ argument of type
{\tt NLPT\_problem\_type}
(see Section~\ref{typeprob}).
It is used to hold data about the problem being solved.
For a new problem, the user must allocate all the array components,
and set values for {\tt nlp\%n}
and the required non-real components of
{\tt nlp\%Jr} depending on what level
of derivatives are required by the model requested. Specifically,
if {\tt control\%jacobian\_available = 2},
all appropriate components of {\tt nlp\%Jr} should be allocated and, 
with the exception of {\tt nlp\%Jr\%val},
filled with data,
Users are free to choose whichever
of the matrix formats described in Section~\ref{galmatrix}
is appropriate for $\bmJ_r$ for their application.

\noindent
The component {\tt nlp\%X} must be set to an initial estimate, $\bmx^{0}$,
of the minimization variables. A good choice will increase the speed
of the package, but the underlying method is designed to converge (at least
to a local solution) from an arbitrary initial guess.
The component {\tt nlp\%COHORT} needs only be allocated and set if there are 
multiple non-overlapping regular simplex constraints \req{c}, while 
the same is true for {\tt nlp\%W} when non-unit weights are required.

\noindent
On exit, the component {\tt nlp\%X} will contain the best estimates of the
minimization variables $\bmx$.

\itt{W} is an \optional\ rank-one array of type \realdp\
whose components specifies the diagonal weighting matrix $\bmW$ that
defines the objective function $f(\bmx)$. If {\tt W} is present, it
must be of length at least $n$, and {\tt W}$(i)$ should contain $w_i > 0$,
$i = 1, \ldots, n$. If $\bmW$ is absent, the weights $w_i$ will all be
taken to be $1.0$.

\noindent
\restrictions {\tt nlp\%n} $> 0$, {\tt nlp\%m\_r} $> 0$ and
{\tt nlp\%Jr\%type} $\in \{
  \mbox{\tt 'DENSE'}, \mbox{\tt 'COORDINATE'}, \mbox{\tt 'SPARSE\_BY\_ROWS'}
 \}$ 
%and
%{\tt nlp\%H\%type} $\in \{
%  \mbox{\tt 'DENSE'}, \mbox{\tt 'COORDINATE'}, \mbox{\tt 'SPARSE\_BY\_ROWS'},
%  \mbox{\tt 'DIAGONAL'} \}$.

\itt{control} is a scalar \intentin\ argument of type
{\tt \packagename\_control\_type}
(see Section~\ref{typecontrol}). Default values may be assigned by calling
{\tt \packagename\_initialize} prior to the first call to
{\tt \packagename\_solve}. 

\itt{inform} is a scalar \intentinout\ argument of type
{\tt \packagename\_inform\_type}
(see Section~\ref{typeinform}).
On initial entry, the  component {\tt inform\%status} must be set to the
value 1. Other entries need note be set.
A successful call to
{\tt \packagename\_solve}
is indicated when the  component {\tt inform\%status} has the value 0.
For other return values of {\tt inform\%status}, see
Sections~\ref{reverse} and \ref{galerrors}.

\itt{data} is a scalar \intentinout\ argument of type
{\tt \packagename\_data\_type}
(see Section~\ref{typedata}). It is used to hold data about the problem being
solved. With the possible exceptions of the components
{\tt data\%eval\_status} and {\tt data\%U} (see Section~\ref{reverse}),
it must not have been altered {\bf by the user} since the last call to
{\tt \packagename\_initialize}.

\itt{userdata} is a scalar \intentinout\ argument of type
{\tt USERDATA\_type} whose components may be used
to communicate user-supplied data to and from the \optional\ subroutines
{\tt eval\_R}, {\tt eval\_Jr}, {\tt eval\_H}, {\tt eval\_JPROD},
{\tt eval\_HPROD}, {\tt eval\_HPRODS} and {\tt eval\_SCALE}
(see Section~\ref{typeuserdata}).

\itt{userdata} is a scalar \intentinout\ argument of type
{\tt USERDATA\_type} whose components may be used
to communicate user-supplied data (see Section~\ref{typeuserdata})
to and from the \optional\ subroutines
 {\tt eval\_R}, {\tt eval\_Jr},
 {\tt eval\_Jr\_prod}, {\tt eval\_Jr\_scol} and {\tt eval\_Jr\_sprod} 
(see below).

\itt{reverse} is an \optional\ scalar \intentinout\ argument of type
{\tt REVERSE\_type}
(see Section~\ref{typereverse}).
It is used to communicate reverse-communication data between the
subroutine and calling program. It is required if either 
(i) {\tt control\%jacobian\_available = 2} and {\tt eval\_Jr} 
(see below) is absent or 
(ii) {\tt control\%jacobian\_available = 1} and 
{\tt eval\_Jr\_prod} or {\tt eval\_Jr\_scol} are absent.
In these cases, the user should monitor {\tt inform\%status} on exit
(see Section~\ref{reverse}).

\itt{eval\_R} is an \optional\
user-supplied subroutine whose purpose is to evaluate the value of the
residual function $\bmr(\bmx)$ at a given vector $\bmx$.
See Section~\ref{rfv} for details.
%but first check Table~\ref{tab-er} to see if the subroutine is not needed.
If {\tt eval\_R} is present,
it must be declared {\tt EXTERNAL} in the calling program.
If {\tt eval\_R} is absent, \solver\ will use reverse communication to
obtain objective function values (see Section~\ref{reverse}).

\itt{eval\_Jr} is an \optional\
user-supplied subroutine whose purpose is to evaluate the value of the
Jacobian of the residual function $\bmJ_r(\bmx)$ at a given vector $\bmx$.
See Section~\ref{jfv} for details.
If {\tt eval\_J} is present,
it must be declared {\tt EXTERNAL} in the calling program.
If {\tt eval\_J} is absent, \solver\ will use reverse communication to
obtain Jacobian values (see Section~\ref{reverse}).

\itt{eval\_Jr\_prod} is an \optional\
user-supplied subroutine whose purpose is to evaluate the value of the product
$\bmp = \bmJ_r(\bmx)\bmv$ or $\bmp = \bmJ_r^T(\bmx)\bmv$ involving the
Jacobian of the residual function $\bmJ_r(\bmx)$, or its transpose, and 
a given vector $\bmv$. See Section~\ref{jvp} for details.
If {\tt eval\_Jr\_prod} is present, it must be declared {\tt EXTERNAL} in the 
calling program. If {\tt eval\_Jr\_prod} is absent, \solver\ will use 
reverse communication to obtain Jacobian-vector products if required 
(see Section~\ref{reverse}).

\itt{eval\_Jr\_scol} is an \optional\
user-supplied subroutine whose purpose is to evaluate the nonzero
components of a specified column of  $\bmJ_r(\bmx)$.
See Section~\ref{jcol} for details,
If {\tt eval\_Jr\_scol} is present,
it must be declared {\tt EXTERNAL} in the calling program.
If {\tt eval\_Jr\_scol} is absent, \solver\ will use reverse communication to
obtain the required column (see Section~\ref{reverse}).

\itt{eval\_Jr\_sprod} is an \optional\
user-supplied subroutine whose purpose is to evaluate the value of the product
$\bmp = \bmJ_r(\bmx)\bmv$ for a given sparse vector $\bmv$, or specified
components of the product $\bmp = \bmJ_r^T(\bmx)\bmv$. involving the
Jacobian of the residual function $\bmJ_r(\bmx)$.
See Section~\ref{sjvp} for details.
If {\tt eval\_Jr\_sprod} is present, it must be declared {\tt EXTERNAL} in the 
calling program. If {\tt eval\_Jr\_sprod} is absent, \solver\ will use 
reverse communication to obtain sparse Jacobian-vector products if required 
(see Section~\ref{reverse}).

\end{description}

%%%%%%% termination subroutine %%%%%%

\subsubsection{The  termination subroutine}
All previously allocated arrays are deallocated as follows:
\vspace*{1mm}

\hspace{8mm}
{\tt CALL \packagename\_terminate( data, control, inform )}

\vspace*{-1mm}
\begin{description}

\itt{data} is a scalar \intentinout\ argument of type
{\tt \packagename\_data\_type}
exactly as for
{\tt \packagename\_solve},
which must not have been altered {\bf by the user} since the last call to
{\tt \packagename\_initialize}.
On exit, array components will have been deallocated.

\itt{control} is a scalar \intentin\ argument of type
{\tt \packagename\_control\_type}
exactly as for
{\tt \packagename\_solve}.

\itt{inform} is a scalar \intentout\ argument of type
{\tt \packagename\_inform\_type}
exactly as for
{\tt \packagename\_solve}.
The component {\tt inform\%status} will be set on exit, and a
successful call to
{\tt \packagename\_terminate}
is indicated when this  component has the value 0.
For other return values of {\tt inform\%status}, see Section~\ref{galerrors}.

\end{description}

%%%%%%%%%%%%%%%%%%%% Function and derivative values %%%%%%%%%%%%%%%%%%%%%

\subsection{Function and derivative values\label{fdv}}

%%%%%%%%%%%%%%%%%%%%%% Residual function value %%%%%%%%%%%%%%%%%%%%%%

\subsubsection{The residual value via internal evaluation\label{rfv}}

If the argument {\tt eval\_R} is present when calling \solver, the
user is expected to provide a subroutine of that name to evaluate the
value of the residual functions $\bmr(\bmx)$.
The routine must be specified as

\def\baselinestretch{0.8}
{\tt
\begin{verbatim}
      SUBROUTINE eval_R( status, X, userdata, R )
\end{verbatim}
}
\def\baselinestretch{1.0}
\noindent whose arguments are as follows:

\begin{description}
\itt{status} is a scalar \intentout\ argument of type \integer,
that should be set to 0 if the routine has been able to evaluate
the residual functions
and to a non-zero value if the evaluation has not been possible.

\ittf{X} is a rank-one \intentin\ array argument of type \realdp\
whose components contain the vector $\bmx$.

\itt{userdata} is a scalar \intentinout\ argument of type
{\tt USERDATA\_type} whose components may be used
to communicate user-supplied data to and from the subroutine
(see Section~\ref{typeuserdata}).

\ittf{R} is a rank-one \intentout\ array argument of type \realdp\
that should be set to the value of the residuals $\bmr(\bmx)$
evaluated at the vector $\bmx$ input in {\tt X}.

\end{description}

%%%%%%%%%%%%%%%%%%%%%% Jacobian values %%%%%%%%%%%%%%%%%%%%%%

\subsubsection{Jacobian values via internal evaluation\label{jfv}}

If the argument {\tt eval\_Jr} is present when calling \solver, the
user is expected to provide a subroutine of that name to evaluate the
values of the residual Jacobian $\bmJ_r(\bmx)$.
The routine must be specified as

\def\baselinestretch{0.8}
{\tt
\begin{verbatim}
       SUBROUTINE eval_Jr( status, X, userdata, Jr_val )
\end{verbatim} }
\def\baselinestretch{1.0}
\noindent whose arguments are as follows:

\begin{description}
\itt{status} is a scalar \intentout\ argument of type \integer,
that should be set to 0 if the routine has been able to evaluate
the residual Jacobian,
and to a non-zero value if the evaluation has not been possible.

\ittf{X} is a rank-one \intentin\ array argument of type \realdp\
whose components contain the vector $\bmx$.

\itt{userdata} is a scalar \intentinout\ argument of type
{\tt USERDATA\_type} whose components may be used
to communicate user-supplied data to and from the subroutine
(see Section~\ref{typeuserdata}).

\itt{Jr\_val} is a scalar \intentout\ argument of type \realdp,
whose components should be set to the values
of the Jacobian $\bmJ_r(\bmx)$
evaluated at the vector $\bmx$ input in {\tt X}.
The values should
be input in the same order as that in which the array indices were
given in {\tt nlp\%J}.

\end{description}

%%%%%%%%%%%%%%%%%%%%%% Jacobian-vector products %%%%%%%%%%%%%%%%%%%%%%

\subsubsection{Jacobian-vector products via internal evaluation\label{jvp}}

If the argument {\tt eval\_Jr\_prod} is present when calling \solver, the
user is expected to provide a subroutine of that name to evaluate the
product $\bmp = \bmJ_r(\bmx) \bmv$ or $\bmp = \bmJ_r^T(\bmx) \bmv$ involving the
product of the residual Jacobian $\bmJ_r(\bmx)$ or its transpose 
$\bmJ_r^T(\bmx)$. The routine must be specified as

\def\baselinestretch{0.8}
{\tt
\begin{verbatim}
       SUBROUTINE eval_Jr_prod( status, X, userdata, transpose, V, P, got_Jr )
\end{verbatim} }
\def\baselinestretch{1.0}
\noindent whose arguments are as follows:

\begin{description}
\itt{status} is a scalar \intentout\ argument of type \integer,
that should be set to 0 if the routine has been able to evaluate the
product $\bmp = \bmJ_r(\bmx) \bmv$ or $\bmp = \bmJ_r^T(\bmx) \bmv$
and to a non-zero value if the evaluation has not been possible.

\ittf{X} is a rank-one \intentin\ array argument of type \realdp\
whose components contain the vector $\bmx$.

\itt{userdata} is a scalar \intentinout\ argument of type
{\tt USERDATA\_type} whose components may be used
to communicate user-supplied data to and from the subroutine
(see Section~\ref{typeuserdata}).

\itt{transpose} is a scalar \intentin\ array argument of type
default that will be set \true\ if the product involves the transpose
of the Jacobian $\bmJ_r^T(\bmx)$ and \false\ if the product involves
the Jacobian $\bmJ_r(\bmx)$ itself.

\ittf{V} is a rank-one \intentin\ array argument of type \realdp\
whose components contain the vector $\bmv$.

\ittf{P} is a rank-one \intentinout\ array argument of type \realdp\
whose components on input contain the vector $\bmp$ and on output the
product $\bmp = \bmJ_r(\bmx) \bmv$ when {\tt \%transpose} is \false\ or
$\bmp = \bmJ_r^T(\bmx) \bmv$ when {\tt \%transpose} is \true.

\itt{got\_Jr} is a scalar \intentin\ array argument of type
default that will be set \true\ if the user has already computed
some information about $\bmJ_r(\bmx)$ at the given {\tt X}, and thus
may avoid unnecessary re-computation. It will be \false\ otherwise.

\end{description}


%%%%%%%%%%%%%%%%%%%%%% Jacobian-vector products %%%%%%%%%%%%%%%%%%%%%%

\subsubsection{Columns of the Jacobian via internal evaluation\label{jcol}}

If the argument {\tt eval\_Jr\_scol} is present when calling \solver, the
user is expected to provide a subroutine of that name to evaluate a specified 
column of the Jacobian $\bmJ_r(\bmx)$. The routine must be specified as

\def\baselinestretch{0.8}
{\tt \begin{verbatim}
       SUBROUTINE eval_Jr_scol( status, X, userdata, index, VAL, ROW, NZ, got_Jr )
\end{verbatim} }
\def\baselinestretch{1.0}
\noindent whose arguments are as follows:

\begin{description}
\itt{status} is a scalar \intentout\ argument of type \integer,
that should be set to 0 if the routine has been able to evaluate the desired 
columns of $\bmJ_r(\bmx)$ and to a non-zero value if the evaluation has not 
been possible.

\ittf{X} is a rank-one \intentin\ array argument of type \realdp\
whose components contain the vector $\bmx$.

\itt{userdata} is a scalar \intentinout\ argument of type
{\tt USERDATA\_type} whose components may be used
to communicate user-supplied data to and from the
subroutine (see Section~\ref{typeuserdata}).

\itt{index} is a scalar \intentout\ argument of type \integer,
that should be set to the index of the required column of $\bmJ_r(\bmx)$.

\ittf{VAL} is a rank-one \intentout\ array argument of type \realdp\
whose first {\tt nz} components on output contain the 
nonzeros of the {\tt index}-th column of $\bmJ_r(\bmx)$.

\itt{ROW} is an \intentinout\ rank-one array of dimension at least
{\tt nz} and type \integer, that must be set to record the list
of indices of nonzero components of the {\tt index}-th column of $\bmJ_r(\bmx)$.

\itt{nz} is an \intentinout\ scalar variable of type \integer,
that must be set to record the number of non-zeros in the {\tt col}-th column
of $\bmJ_r(\bmx)$.

\itt{got\_Jr} is a scalar \intentin\ array argument of type
default that will be set \true\ if the user has already computed
some information about $\bmJ_r(\bmx)$ at the given {\tt X}, and thus
may avoid unnecessary re-computation. It will be \false\ otherwise.

\end{description}

%%%%%%%%%%%%%%%%%%%%%% Jacobian-vector products %%%%%%%%%%%%%%%%%%%%%%

\subsubsection{Jacobian-vector sub-products via internal evaluation\label{sjvp}}

If the argument {\tt eval\_Jr\_sprod} is present when calling \solver, the
user is expected to provide a subroutine of that name to evaluate the
product of the Jacobian $\bmJ_r(\bmx)$, or its transpose, with a given 
vector $\bmv$.
Here, either only a subset of the components of the vector $\bmv$ are nonzero,
or only a subset of the components product $\bmJ_r^T(\bmx) \bmv$ are required.
The routine must be specified as

\def\baselinestretch{0.8}
{\tt \begin{verbatim}
       SUBROUTINE eval_Jr_sprod( status, X, userdata, transpose, V, P, FREE, n_free, got_Jr )
\end{verbatim} }
\def\baselinestretch{1.0}
\noindent whose arguments are as follows:

\begin{description}
\itt{status} is a scalar \intentout\ argument of type \integer,
that should be set to 0 if the routine has been able to evaluate the
sparse product $\bmJ_r \bmv$ or $\bmJ_r^T\bmv$
and to a non-zero value if the evaluation has not been possible.

\ittf{X} is a rank-one \intentin\ array argument of type \realdp\
whose components contain the vector $\bmx$.

\itt{userdata} is a scalar \intentinout\ argument of type
{\tt USERDATA\_type} whose components may be used
to communicate user-supplied data to and from the subroutine
(see Section~\ref{typeuserdata}).

\itt{transpose} is a scalar \intentin\ array argument of type
default that will be set \true\ if the product involves the transpose
of the Jacobian $\bmJ_r^T(\bmx)$ and \false\ if the product involves
the Jacobian $\bmJ_r(\bmx)$ itself.

\ittf{V} is a rank-one \intentin\ array argument of type \realdp\
whose components contain the vector $\bmv$. If {\tt transpose} is \false\
only those components whose indices {\tt FREE(:n\_free)} (see below)
will by set, and the remainder will be presumed to be zero.

\ittf{P} is a rank-one \intentout\ array argument of type \realdp\
whose components on output must be set to the
product $\bmJ_r(\bmx) \bmv$ when {\tt \%transpose} is \false. If
{\tt \%transpose} is \true, components with indices {\tt FREE(:n\_free)}
(see below) should be set to the corresponding components of the
product $\bmJ_r^T \bmv$, and the remaining components ignored.

\itt{FREE} is a rank-one \intentin\ array argument of
type \integer\ that flags the input components of $\bmv$ that are set
(when {\tt transpose} is \false) or output components of $\bmJ_r^T(\bmx) \bmv$ 
that are required (when {\tt transpose} is \true). Specifically, only indices
{\tt FREE(:n\_free)} of the relevant vector is set or required, and the
remainder should be treated as zero (if {\tt transpose} is \false)
or ignored  (if {\tt transpose} is \true).

\itt{n\_free} is a scalar \intentin\ argument of type \integer,
that specifies the number of components of {\tt FREE} that need be considered.

\itt{got\_Jr} is a scalar \intentin\ array argument of type
default that will be set \true\ if the user has already computed
some information about $\bmJ_r(\bmx)$ at the given {\tt X}, and thus
may avoid unnecessary re-computation. It will be \false\ otherwise.

\end{description}

%%%%%%%%%%%%%%%%%%%%%% Reverse communication %%%%%%%%%%%%%%%%%%%%%%%%

\subsection{\label{reverse}Reverse Communication Information}

A positive value of {\tt inform\%status} on exit from
{\tt \packagename\_solve}
indicates that
\solver\ is seeking further information---this will happen
if the user has chosen not to evaluate function or
derivative values internally (see Section~\ref{fdv}).
The user should compute the required information and re-enter \solver\
with {\tt inform\%status} and all other arguments (except those specifically
mentioned below) unchanged.

Possible values of {\tt inform\%status} and the information required are
\begin{description}
\ittf{2.} The user should compute the vector of residual functions
     $\bmr(\bmx)$ at the point $\bmx$ indicated in {\tt nlp\%X}.
     The required value should be set in {\tt nlp\%R}, and
     {\tt data\%eval\_status} should be set to 0. If the user is
     unable to evaluate $\bmr(\bmx)$---for instance, if one or more of the
     residual functions is undefined at $\bmx$---the user need not set
     {\tt nlp\%R}, but should then set {\tt data\%eval\_status}
     to a non-zero value.

\ittf{3.} The user should compute the Jacobian matrix $\bmJ_r(\bmx)$ of the
     residuals $\bmr(\bmx)$ at the point $\bmx$ indicated in {\tt nlp\%X}.
     The $l$-th component of the Jacobian stored according to the scheme
     input in the remainder of {\tt nlp\%Jr} (see Section~\ref{typeprob})
     should be set in {\tt nlp\%Jr\%val(l)}, for $l = 1, \ldots,$
     {\tt  nlp\%Jr\%ne} and {\tt data\%eval\_status} should be set to 0.
     If the user is unable to  evaluate a component of $\bmJ_r(\bmx)$---for
     instance, if a component of the Jacobian is
     undefined at $\bmx$---the user need not set {\tt nlp\%Jr\%val}, but
     should then set {\tt data\%eval\_status} to a non-zero value.

\ittf{4.} The user should compute the product $\bmp = \bmJ_r(\bmx) \bmv$
     involving the residual Jacobian $\bmJ_r(\bmx)$ at the point $\bmx$, given
     in {\tt nlp\%X}, and a given vector $\bmv$.
     The vector $\bmv$ is provided in {\tt reverse\%V},
     the resulting $\bmp$ should be placed in {\tt reverse\%P},
     and  {\tt reverse\%eval\_status} should be set to 0. If the user is
     unable to evaluate the product---for instance, if a component of the 
     Jacobian is undefined at $\bmx$---the user need not set {\tt reverse\%P}, 
     but should then set {\tt reverse\%eval\_status} to a non-zero value.

\ittf{5.} The user should compute the product $\bmp = \bmJ_r^T(\bmx) \bmv$ 
     involving the transpose of the residual Jacobian $\bmJ_r(\bmx)$ 
     at the point $\bmx$, given in {\tt nlp\%X}, and a given vector $\bmv$.
     The vector $\bmv$ is provided in {\tt reverse\%V},
     the resulting $\bmp$ should be placed in {\tt reverse\%P},
     and  {\tt reverse\%eval\_status} should be set to 0. If the user is
     unable to evaluate the product---for instance, if a component of the 
     Jacobian is undefined at $\bmx$---the user need not set {\tt reverse\%P}, 
     but should then set {\tt reverse\%eval\_status} to a non-zero value.


\ittf{6.} The user should compute the $j$-th column of $\bmJ_r(\bmx)$, with
     $j$ provided in {\tt reverse\%index}, at the point $\bmx$ given in 
     {\tt nlp\%X}. The resulting {\em nonzeros} and their corresponding 
     row indices of the $j$-th column of $\bmJ_r(\bmx)$ must be placed in 
     {\tt reverse\%P(1 : reverse\%lp)} and 
     {\tt reverse\%IP(1 : reverse\%lp)} with {\tt reverse\%lp} set accordingly,
     and {\tt reverse\%eval\_status} should be set to 0. If the user is
     unable to evaluate the required column---for instance, if a component of 
     the Jacobian is undefined at $\bmx$---the user need not set 
     {\tt reverse\%P}, {\tt reverse\%IP} and {\tt reverse\%nz},
     but should then set {\tt reverse\%eval\_status} to a non-zero value.

\ittf{7.} The user should compute the product $\bmp = \bmJ_r(\bmx) \bmv$
     involving the residual Jacobian $\bmJ_r(\bmx)$ at the point $\bmx$, given
     in {\tt nlp\%X}, and a given sparse vector $\bmv$, whose nonzeros
     are in positions {\tt reverse\%iv(reverse\%lvl:reverse\%lvu)}
     of {\tt reverse\%V}. The resulting $\bmp$ should be placed in 
     {\tt reverse\%P} and {\tt reverse\%eval\_status} should be set to 0. 
     If the user is unable to evaluate the product---for instance, if a 
     component of the Jacobian is undefined at $\bmx$---the user need not 
     set {\tt reverse\%P}, but should then set {\tt reverse\%eval\_status} 
     to a non-zero value.

\ittf{8.} The user should compute selected components of the product 
     $\bmp = \bmJ_r^T(\bmx) \bmv$ involving the transpose of the residual 
     Jacobian $\bmJ_r(\bmx)$ at the point $\bmx$, given
     in {\tt nlp\%X}, and a given vector $\bmv$. Only components \linebreak
     {\tt reverse\%IV(reverse\%lvl:reverse\%lvu)} of $\bmp$ should
     be computed, and recorded in \linebreak
     {\tt reverse\%P(reverse\%IV(reverse\%lvl:reverse\%lvu))}, and
     {\tt reverse\%eval\_status} should be set to 0. 
     If the user is unable to evaluate the product---for instance, if a 
     component of the Jacobian is undefined at $\bmx$---the user need not 
     set {\tt reverse\%P}, but should then set {\tt reverse\%eval\_status} 
     to a non-zero value.

\ittf{9.} The user has the opportunity to replace the estimate $\bmx$ in 
     {\tt nlp\%X} by an improved value $\bmx^+$ for which 
     $f(\bmx^+) \leq f(\bmx)$; in that case {\tt nlp\%R} must also be reset to
     hold $\bmr(\bmx^+)$.

\end{description}


%%%%%%%%%%%%%%%%%%%%%% Warning and error messages %%%%%%%%%%%%%%%%%%%%%%%%

\galerrors
A negative value of {\tt inform\%status} on exit from
{\tt \packagename\_solve}
or
{\tt \packagename\_terminate}
indicates that an error has occurred. No further calls should be made
until the error has been corrected. Possible values are:

\begin{description}

\itt{\galerrallocate.} An allocation error occurred.
A message indicating the offending
array is written on unit {\tt control\%error}, and the returned allocation
status and a string containing the name of the offending array
are held in {\tt inform\%alloc\_\-status}
and {\tt inform\%bad\_alloc}, respectively.

\itt{\galerrdeallocate.} A deallocation error occurred.
A message indicating the offending
array is written on unit {\tt control\%error} and the returned allocation
status and a string containing the name of the offending array
are held in {\tt inform\%alloc\_\-status}
and {\tt inform\%bad\_alloc}, respectively.

\itt{\galerrrestrictions.}
  The restriction {\tt nlp\%n} $> 0$ and {\tt nlp\%m\_r} $> 0$,
  or requirements that
  {\tt nlp\%Jr\%type}  contains its relevant string
  \mbox{\tt 'DENSE'}, \mbox{\tt 'COORDINATE'} 
  has been violated.

%\itt{-3.} At least one of the arrays
% {\tt p\%A\_val}, {\tt p\%A\_row}, {\tt p\%A\_col},
% {\tt p\%H\_val}, {\tt p\%H\_row} or {\tt p\%H\_col},
% is not large enough to hold the original, or reordered, matrices $\bmJ_r$
% or $\bmH$.

%\itt{\galerrunbounded.}  The objective function appears to be unbounded
% from below on the feasible set.

\itt{\galerranalysis.} The analysis phase of the factorization failed;
 the return status from the factorization
    package is given in the component {\tt inform\%fac\-t\-or\_status}.

\itt{\galerrfactorization.} The factorization failed;
 the return status from the factorization
    package is given in the component {\tt inform\%fac\-t\-or\_status}.

\itt{\galerrsolve.} The solution of a set of linear equations
 using factors from the factorization package failed;
 the return status from the factorization
    package is given in the component {\tt inform\%fac\-t\-or\_status}.

\itt{\galerrillconditioned.} The problem is so ill-conditioned that
  further progress is impossible.

\itt{\galerrtinystep.} The step is too small to make further impact.

\itt{\galerrmaxiterations.} Too many iterations have been performed.
  This may happen if {\tt control\%maxit} is too small, but may also 
  be symptomatic of a badly scaled problem.

\itt{\galerrcpulimit.} The elapsed CPU or system clock time limit has been
    reached. This may happen if either {\tt control\%cpu\_time\-\_limit} or
    {\tt control\%clock\_time\_limit} is too small, but may also be symptomatic
    of a badly scaled problem.

\itt{\galerralive.} The user has forced termination of \solver\
     by removing the file named {\tt control\%a\-live\_file} from unit
     unit {\tt control\%alive\_unit}.

\end{description}

%%%%%%%%%%%%%%%%%%%%%% Further features %%%%%%%%%%%%%%%%%%%%%%%%

\galfeatures
\noindent In this section, we describe an alternative means of setting
control parameters, that is components of the variable {\tt control} of type
{\tt \packagename\_control\_type}
(see Section~\ref{typecontrol}),
by reading an appropriate data specification file using the
subroutine {\tt \packagename\_read\_specfile}. This facility
is useful as it allows a user to change  {\tt \packagename} control parameters
without editing and recompiling programs that call {\tt \packagename}.

A specification file, or specfile, is a data file containing a number of
"specification commands". Each command occurs on a separate line,
and comprises a "keyword",
which is a string (in a close-to-natural language) used to identify a
control parameter, and
an (optional) "value", which defines the value to be assigned to the given
control parameter. All keywords and values are case insensitive,
keywords may be preceded by one or more blanks but
values must not contain blanks, and
each value must be separated from its keyword by at least one blank.
Values must not contain more than 30 characters, and
each line of the specfile is limited to 80 characters,
including the blanks separating keyword and value.

The portion of the specification file used by
{\tt \packagename\_read\_specfile}
must start
with a "{\tt BEGIN \packagename}" command and end with an
"{\tt END}" command.  The syntax of the specfile is thus defined as follows:
\begin{verbatim}
  ( .. lines ignored by SNLS_read_specfile .. )
    BEGIN SNLS
       keyword    value
       .......    .....
       keyword    value
    END
  ( .. lines ignored by SNLS_read_specfile .. )
\end{verbatim}
where keyword and value are two strings separated by (at least) one blank.
The ``{\tt BEGIN \packagename}'' and ``{\tt END}'' delimiter command lines
may contain additional (trailing) strings so long as such strings are
separated by one or more blanks, so that lines such as
\begin{verbatim}
    BEGIN SNLS SPECIFICATION
\end{verbatim}
and
\begin{verbatim}
    END SNLS SPECIFICATION
\end{verbatim}
are acceptable. Furthermore,
between the
``{\tt BEGIN \packagename}'' and ``{\tt END}'' delimiters,
specification commands may occur in any order.  Blank lines and
lines whose first non-blank character is {\tt !} or {\tt *} are ignored.
The content
of a line after a {\tt !} or {\tt *} character is also
ignored (as is the {\tt !} or {\tt *}
character itself). This provides an easy manner to "comment out" some
specification commands, or to comment specific values
of certain control parameters.

The value of a control parameters may be of three different types, namely
integer, logical or real.
Integer and real values may be expressed in any relevant Fortran integer and
floating-point formats (respectively). Permitted values for logical
parameters are "{\tt ON}", "{\tt TRUE}", "{\tt .TRUE.}", "{\tt T}",
"{\tt YES}", "{\tt Y}", or "{\tt OFF}", "{\tt NO}",
"{\tt N}", "{\tt FALSE}", "{\tt .FALSE.}" and "{\tt F}".
Empty values are also allowed for
logical control parameters, and are interpreted as "{\tt TRUE}".

The specification file must be open for
input when {\tt \packagename\_read\_specfile}
is called, and the associated device number
passed to the routine in device (see below).
Note that the corresponding
file is {\tt REWIND}ed, which makes it possible to combine the specifications
for more than one program/routine.  For the same reason, the file is not
closed by {\tt \packagename\_read\_specfile}.

\subsubsection{To read control parameters from a specification file}
\label{readspec}

Control parameters may be read from a file as follows:
\hskip0.5in
\def\baselinestretch{0.8} {\tt
\begin{verbatim}
     CALL SNLS_read_specfile( control, device )
\end{verbatim}
}
\def\baselinestretch{1.0}

\begin{description}
\itt{control} is a scalar \intentinout\ argument of type
{\tt \packagename\_control\_type}
(see Section~\ref{typecontrol}).
Default values should have already been set, perhaps by calling
{\tt \packagename\_initialize}.
On exit, individual components of {\tt control} and
{\tt control\%subproblem\_control}
may have been changed
according to the commands found in the specfile. Specfile commands and
the component (see Section~\ref{typecontrol}) of {\tt control}
and {\tt control\%subproblem\_control}
that each affects are given in Table~\ref{specfile}.

\bctable{|l|l|l|}
\hline
  command & component of {\tt control} & value type \\
\hline
  {\tt error-printout-device} & {\tt \%error} & integer \\
  {\tt printout-device} & {\tt \%out} & integer \\
  {\tt print-level} & {\tt \%print\_level} & integer \\
  {\tt start-print} & {\tt \%start\_print} & integer \\
  {\tt stop-print} & {\tt \%stop\_print} & integer \\
  {\tt iterations-between-printing} & {\tt \%print\_gap} & integer \\
  {\tt maximum-number-of-iterations} & {\tt \%maxit} & integer \\
  {\tt alive-device} & {\tt \%alive\_unit} & integer \\
  {\tt jacobian-available} & {\tt \%jacobian\_available} & integer \\
  {\tt subproblem-solver} & {\tt \%subproblem\_solver} & integer \\
  {\tt history-length-for-non-monotone-descent} & {\tt \%non\_monotone} & integer \\
  {\tt weight-update-strategy} & {\tt weight\_update\_strategy} & integer \\
  {\tt absolute-residual-accuracy-required} & {\tt \%stop\_c\_absolute} & real \\
  {\tt relative-residual-reduction-required} & {\tt \%stop\_c\_relative} & real \\
  {\tt absolute-gradient-accuracy-required} & {\tt \%stop\_pg\_absolute} & real \\
  {\tt relative-gradient-reduction-required} & {\tt \%stop\_pg\_relative} & real \\
  {\tt minimum-relative-step-allowed} & {\tt \%stop\_s} & real \\
  {\tt relative-gradient-switch-tolerance} & {\tt \%stop\_pg\_switch} & real \\
  {\tt initial-regularization-weight} & {\tt \%initial\_weight} & real \\
  {\tt minimum-regularization-weight} & {\tt \%minimum\_weight} & real \\
  {\tt successful-iteration-tolerance} & {\tt \%eta\_successful} & real \\
  {\tt very-successful-iteration-tolerance} & {\tt \%eta\_very\_successful} & real \\
  {\tt too-successful-iteration-tolerance} & {\tt \%eta\_too\_successful} & real \\
  {\tt regularization-weight-minimum-decrease-factor} & {\tt \%weight\_decrease\_min} & real \\
  {\tt regularization-weight-decrease-factor} & {\tt \%weight\_decrease} & real \\
  {\tt regularization-weight-increase-factor} & {\tt \%weight\_increase} & real \\
  {\tt regularization-weight-maximum-increase-factor} & {\tt \%weight\_increase\_max} & real \\
  {\tt maximum-cpu-time-limit} & {\tt \%cpu\_time\_limit} & real \\
  {\tt maximum-clock-time-limit} & {\tt \%clock\_time\_limit} & real \\
  {\tt try-newton-acceleration}  & {\tt \%newton\_acceleration} & logical \\
  {\tt choose-magic-step}  & {\tt \%magic\_step} & logical \\
  {\tt print-objective}  & {\tt \%print\_obj} & logical \\
  {\tt space-critical}   & {\tt \%space\_critical} & logical \\
  {\tt deallocate-error-fatal}   & {\tt \%deallocate\_error\_fatal} & logical \\
  {\tt alive-filename} & {\tt \%alive\_file} & character \\
\hline

\ectable{\label{specfile}Specfile commands and associated
components of {\tt control}.}

\itt{device} is a scalar \intentin argument of type \integer,
that must be set to the unit number on which the specfile
has been opened. If {\tt device} is not open, {\tt control} will
not be altered and execution will continue, but an error message
will be printed on unit {\tt control\%error}.

\end{description}

%%%%%%%%%%%%%%%%%%%%%% Information printed %%%%%%%%%%%%%%%%%%%%%%%%

\galinfo
If {\tt control\%print\_level} is positive, information about the progress
of the algorithm will be printed on unit {\tt control\-\%out}.
If {\tt control\%print\_level} $= 1$, a single line of output will be
produced for each iteration of the process.
This will include the values of the objective function and the norm of its
projected gradient, the ratio of actual to predicted decrease following the 
step, the
value of the regularization weight and the time taken so far. 

If {\tt control\%print\_level} $\geq 2$ this
output will be increased to provide significant detail of each iteration.
This extra output includes residuals of the linear systems solved, and,
for larger values of {\tt control\%print\_level}, values of the variables
and gradients. Further details concerning the attempted solution the model
may be obtained by increasing
{\tt control\%SLLS\_control\%print\_level}
and
{\tt control\%SLLSB\_control\%print\_level}.
See the specification sheets for the packages
{\tt \libraryname\_SLLS} and {\tt \libraryname\_SLLSB}
for details.

%%%%%%%%%%%%%%%%%%%%%% GENERAL INFORMATION %%%%%%%%%%%%%%%%%%%%%%%%

\galgeneral

\galcommon None.
\galworkspace Provided automatically by the module.
\galroutines None.
\galmodules {\tt \packagename\_solve} calls the \galahad\ packages
{\tt GALAHAD\_CLOCK},
{\tt GALAHAD\_SY\-M\-BOLS},
{\tt GALA\-HAD\-\_NLPT},
{\tt GALAHAD\_USERDATA},
{\tt GALAHAD\_REVERSE},
{\tt GALAHAD\_SPECFILE},
{\tt GALAHAD\_QPT},
{\tt GALAHAD\_SLLS},
{\tt GALA\-HAD\_SLLSB},
{\tt GALAHAD\_\-SPACE},
{\tt GALAHAD\_MOP},
{\tt GALAHAD\_NORMS}
and
{\tt GALAHAD\_STRING}.
\galio Output is under control of the arguments
 {\tt control\%error}, {\tt control\%out} and
{\tt control\%print\_level}.
\galrestrictions {\tt nlp\%n} $> 0$, {\tt nlp\%m\_r} $> 0$,
{\tt nlp\%Jr\%type} $\in \{
  \mbox{\tt 'DENSE'}, \mbox{\tt 'COORDINATE'}, \mbox{\tt 'SPARSE\_BY\_ROWS'}\}$.
\galportability ISO Fortran~95 + TR 15581 or Fortran~2003.
The package is thread-safe.

%%%%%%%%%%%%%%%%%%%%%% METHOD %%%%%%%%%%%%%%%%%%%%%%%%

\galmethod
An adaptive regularization method is used.
In this, an improvement to a current
estimate of the required minimizer, $\bmx_k$ is sought by computing a
step $\bms_k$. The step is chosen to approximately minimize a model $t_k(\bms)$
of $f_{\rho,r}(\bmx_k+\bms)$ that includes a weighted regularization term
$\half \sigma_k \|s\|_2^2$
for some specified positive weight $\sigma_k$. The quality of the
resulting step $\bms_k$ is assessed by computing the "ratio"
%$(f_{\rho,p}(\bmx_k) - f_{\rho,p}(\bmx_k+\bms_k))/(t_k(\bmzero)-t_k(\bms_k))$.
$(f(\bmx_k) - f(\bmx_k + \bms_k))/(t_k(\bmzero) - t_k(\bms_k))$.
The step is deemed to have succeeded if the ratio exceeds a given $\eta_s > 0$,
and in this case $\bmx_{k+1} = \bmx_k + \bms_k$. Otherwise
$\bmx_{k+1} = \bmx_k$, and the weight is increased by powers of a given
increase factor up to a given limit. If the ratio is larger than
$\eta_v \geq \eta_d$, the weight will be decreased by powers of a given
decrease factor again up to a given limit. The method will terminate
as soon as $f(\bmx_k)$ or
$\|P[\bmx_k - \nabla_x f(\bmx_k)]-\bmx_k\|$ is smaller than a specified value.

The step $\bms_k$ may be computed either by employing a projected-gradient
method to minimize the model within the simplex constraint set 
$\calC(\bmx_k+\bms)$ using the \galahad\ module {\tt SLLS}, or by
applying the interior-point method available in the module {\tt SLLSB}
to the same subproblem. Experience has shown that it can be beneficial
to use the latter method during early iterations, but to switch to the 
former as the iterates approach convergence.

\vspace*{1mm}

\galreferences
\vspace*{1mm}

\noindent
Generic regularization methods are described in detail in
\vspace*{1mm}

\noindent
C. Cartis,  N. I. M. Gould and Ph. L. Toint,
``Evaluation complexity of algorithms for nonconvex optimization''
SIAM-MOS Series on Optimization (2022),
\vspace*{1mm}

\noindent
and uses ``tricks'' as suggested in
\vspace*{1mm}

\noindent
N. I. M. Gould, M. Porcelli and Ph. L. Toint,
``Updating the regularization parameter in the adaptive cubic regularization
algorithm''.
{\em Computational Optimization and Applications}
{\bf 53(1)} (2012) 1--22.

\noindent
The specific methods employed here are discussed in
\vspace*{1mm}

\noindent
N. I. M. Gould et.\ al.
``Nonlinear least-squares over unit simplices''.
Rutherford Appleton Laboratory,
Oxfordshire, England (2026) in preparation.

%%%%%%%%%%%%%%%%%%%%%% EXAMPLE %%%%%%%%%%%%%%%%%%%%%%%%

\galexamples
Suppose we have the parametric residual
\disp{\bmr(\bmx) = 
 \vect{ x_1 x_2 - p \\ x_2 x_3 - 1 \\ x_3 x_4 - 1 \\ x_4 x_5 -1}}
and wish to minimize the least-squares objective \req{f}, with unit weights
$\bmw = 1$ and parameter value $p = 4$, over the intersection of the two
regular simplices
\disp{x_1 + x_4 = 1,\; x_1, x_4 \geq 0 \tim{and} 
 x_2 + x_5 = 1,\; x_2, x_5 \geq 0.}
Noting that 
\disp{\bmJ_r(\bmx) = \mat{ccccc}{ x_2 & x_1 \\ & x_3 & x_2 \\ 
& & x_4 & x_3 \\ & & & x_5 & x_4 },}
assigning variables $x_1$ and $x_4$ to cohort 1, $x_2$ and $x_5$ to cohort 2
(and recording that $x_3$ is unconstrained), 
and starting from the initial guess $\bmx = (0.5,0.5,0.5,0.5,0.5)$,
we may use the following code:

{\tt \small
\VerbatimInput{\packageexample}
}
\noindent
Notice how the parameter $p$ is passed to the function evaluation
routines via the {\tt real} component of the derived type {\tt USERDATA\_type}.
The code produces the following output:
{\tt \small
\VerbatimInput{\packageresults}
}
\noindent
If the Jacobian is unavailable, but products of the form
$\bmJ_r(\bmx) \bmv$ and $\bmJ_r^T(\bmx) \bmv$ 
are, the same problem may be solved as follows:

{\tt \small
\VerbatimInput{\packageexampleb}
}
\noindent
This produces the same output.
%following output:
%{\tt \small
%\VerbatimInput{\packageresultsb}
%}
%\noindent

If the user prefers to provide function and Jacobian information
without calls to specified routines, the following code is appropriate:
%Note the product with the user-provided preconditioner
%\disp{
%\bmS(\bmx) = \mat{ccc}{\half & 0 & 0 \\ 0 & \half & 0 \\ 0 & 0 & \quarter }
%}
%which is a suitable approximation to the inverse of the Hessian:

{\tt \small
\VerbatimInput{\packageexamplec}
}
\noindent
This produces the same output as in the previous case.
%This produces the following output:
%{\tt \small
%\VerbatimInput{\packageresultsc}
%}
\noindent


\end{document}
